\documentclass[11pt]{article}

\usepackage[review]{acl}
\usepackage{times}
\usepackage{latexsym}
\usepackage[T1]{fontenc}
\usepackage[utf8]{inputenc}
\usepackage{microtype}
\usepackage{inconsolata}
\usepackage{graphicx}
\usepackage{amsmath}
\usepackage{amssymb}
\usepackage{booktabs}
\usepackage{subcaption}
\usepackage{xcolor}

\newcommand{\ph}[1]{\textcolor{gray}{[#1]}}
\newcommand{\tcrg}{\textsc{TCRG}}
\newcommand{\todo}[1]{\textcolor{red}{[#1]}}

\title{Why Steering Methods Occasionally Fail:\\
A Geometrical Analysis of Task-Conditioned Refusal in Aligned LLMs}

\author{First Author \\
  Affiliation / Address \\
  \texttt{email@domain} \\\And
  Second Author \\
  Affiliation / Address \\
  \texttt{email@domain} \\}

\begin{document}
\maketitle

%% -----------------------------------------------------------
\begin{abstract}
%% -----------------------------------------------------------

Representation steering has become a primary tool for controlling refusal behavior in safety-aligned language models, yet we observe that these interventions are unreliable: some instruction types resist correction, strong interventions produce incoherent outputs, and effectiveness varies widely across contexts. We investigate the mechanistic roots of these failures through causal attribution, vocabulary projection, and directional geometry analysis across six NLP tasks. Our findings suggest that refusal is not always mediated by a single universal circuit but by a family of task-conditioned sub-circuits whose geometric properties vary systematically. This circuit-level heterogeneity directly explains when and why steering succeeds or fails for safety refusals, and suggests principled directions for more reliable interventions.

\end{abstract}

%% -----------------------------------------------------------
\section{Introduction}
%% -----------------------------------------------------------

Safety alignment in large language models (LLMs) introduces refusal as the primary protective response to harmful inputs \citep{ouyang2022instructgpt,bai2022constitutional}. Mechanistic interpretability (MI) has made substantial progress in localizing this behavior: contrastive activation analysis reveals that a compact direction in the residual stream mediates refusal across a wide range of models, and that ablating this direction disables refusal with minimal capability cost \citep{arditi2024refusal,zhao2025refusal}. Building on this, steering methods modify intermediate representations at inference time to control safety behavior---adding or removing directions associated with refusal, truthfulness, or other target properties \citep{zou2023representation,panickssery2023steering,wangf2025refusal}.

In practice, these methods succeed in many settings but fail in others. Over-refusal on benign task-wrapped inputs is a persistent problem: models refuse translation or sentiment analysis requests involving sensitive content at rates far exceeding their refusal of direct harmful inputs, yet steering interventions designed to correct this often have limited effect or degrade output coherence \citep{safeconstellations2025,braun2025unreliability}. The question of \emph{why} a steering vector works for one instruction type but not another remains largely unanswered at the circuit level.

\begin{figure*}[t]
  \centering
  \fbox{
    \begin{minipage}{0.96\linewidth}
      \centering
      \vspace{0.8em}
      \textbf{Figure 1 --- Overview: task-conditioned refusal circuit structure}

      \smallskip
      \emph{Suggested content:} A three-panel overview figure.
      \textbf{Left:} Schematic of a transformer residual stream showing that different task types (sentiment, translation, code) activate different subsets of attention heads (colored arrows), with a shared ``universal'' head cluster and task-specific head clusters.
      \textbf{Center:} UMAP of last-token residual stream states for all six task conditions, colored by task type, with refusal vs. non-refusal samples shown as filled vs. open circles. Tasks cluster distinctly, with translation and sentiment forming tight refusal clusters.
      \textbf{Right:} Bar chart of per-task bypass efficacy for the universal DIM direction (gray) vs. task-specific directions (colored), showing the large gap for translation and sentiment.
      \vspace{0.8em}
    \end{minipage}
  }
  \caption{Refusal in safety-aligned LLMs is mediated by task-conditioned circuits. The universal difference-in-means direction captures the shared component of these circuits but misses task-specialized sub-circuits, explaining why over-refusal persists for certain instruction types regardless of intervention strength.}
  \label{fig:overview}
\end{figure*}

This paper provides a circuit-level account of that question. We analyze refusal circuits across six NLP task conditions using a combination of activation patching, vocabulary projection via the logit lens, and per-task directional geometry measurements. Our analysis reveals three interconnected phenomena: the optimal refusal direction diverges meaningfully across task types; refusal commitment solidifies at task-dependent layers, leaving different-sized windows for effective intervention; and over-used task types whose refusal is mediated by higher-dimensional feature subspaces are both harder to steer and more sensitive to over-steering collapse.

Together these findings explain the steering failures documented across recent literature as consequences of a single underlying structure: task-conditioned circuit heterogeneity. They also suggest concrete, actionable design principles for steering methods that take task context into account.

Our contributions are: (i) the Task-Conditional Refusal Geometry (\tcrg{}) score, a quantitative measure of per-task directional divergence validated across four model families; (ii) a characterization of universal versus task-specialized attention heads and MLP neurons via activation patching; (iii) a pre-failure diagnostic based on early-layer vocabulary projection entropy; and (iv) a unified mechanistic account that connects refusal circuit geometry to the empirical failure patterns of existing steering approaches.

%% -----------------------------------------------------------
\section{Background}
%% -----------------------------------------------------------
\label{sec:background}

\paragraph{Refusal localization.}
A central result in safety MI is that refusal behavior can be localized to a compact direction in the residual stream. \citet{arditi2024refusal} extract this direction via difference-in-means (DIM) over contrastive harmful/harmless activation pairs, showing it is necessary and sufficient for refusal across 13 models. \citet{zhao2025refusal} extend this to reasoning models, identifying refusal subspaces via causal attribution. At the component level, \citet{zhou2025safety} showed that a small subset of attention heads disproportionately mediates safety behavior, and \citet{zhao2025d} identified safety-governing neurons whose selective modification edits safety behavior with minimal side effects. SAE-based analyses \citep{templeton2024scaling,yeo2025sae,goyal2025detox} provide a feature-level decomposition, showing that harm- and refusal-related features are monosemantic and localized.

\paragraph{Steering vectors and their limitations.}
Vector arithmetic on the residual stream---adding or ablating directions associated with target behaviors---is the dominant paradigm for inference-time safety control \citep{zou2023representation,panickssery2023steering}. \citet{wangf2025refusal} propose a training-free ablation that selectively removes false refusal while preserving true refusal. Across these methods, a shared assumption is that the target behavior is encoded as a direction that is approximately constant across input contexts.

\citet{braun2025unreliability} examine when this assumption breaks down, identifying three failure conditions: high steering magnitudes, distributional distance between the extraction and application context, and feature entanglement in the steering vector. \citet{hedstrom2025steer} complement this with a decision-theoretic analysis showing that abstention from steering is optimal when the model lacks stable internal representations for the target behavior. \citet{wang2025adaptive} extend the concern to hallucination: a single universal vector for truthfulness fails across diverse hallucination categories, requiring per-category adaptive vectors.

\paragraph{Task-conditioned representations.}
\citet{safeconstellations2025} show that LLMs encode task identity through stable trajectory ``constellation'' patterns in layer-wise activation space, and that these patterns are more predictive of refusal than input harmfulness. Their inference-time steering method reduces over-refusal by up to 73\% by targeting layers of maximal task-trajectory divergence. This establishes the empirical phenomenon our work explains mechanistically.

\paragraph{Feature geometry.}
The linear representation hypothesis \citep{park2023linear,elhage2021mathematical} holds that concepts are encoded as one-dimensional directions. \citet{engels2025nonlinear} challenge this for a class of features: days of the week and months of the year are encoded as circular manifolds in GPT-2 and Mistral 7B, and require multi-dimensional representations for interventions to succeed. SAE clustering can automatically discover such structures. We apply this geometric lens to refusal features, finding that task-conditioned over-refusal requires substantially higher-dimensional representations than direct harmful-input refusal.

%% -----------------------------------------------------------
\section{Method}
%% -----------------------------------------------------------
\label{sec:method}

\paragraph{Task-Conditional Refusal Geometry.}
Following \citet{arditi2024refusal}, we extract a per-task DIM direction for each task condition $k$:
\begin{equation}
  \hat{\mathbf{r}}^{(k)} = \frac{\boldsymbol{\mu}^{(k)}_\text{harm} - \boldsymbol{\mu}^{(k)}_\text{harmless}}{\|\boldsymbol{\mu}^{(k)}_\text{harm} - \boldsymbol{\mu}^{(k)}_\text{harmless}\|}
\end{equation}
and define the \tcrg{} score as the mean pairwise cosine distance across all task pairs:
\begin{equation}
  \tcrg = \frac{1}{\binom{K}{2}} \sum_{k < k'} \bigl(1 - \hat{\mathbf{r}}^{(k)} \cdot \hat{\mathbf{r}}^{(k')}\bigr)
\end{equation}
A \tcrg{} of zero means all task directions are aligned (single-direction account holds); non-trivial values indicate task-conditioned geometry.

\paragraph{Activation patching.}
We apply path patching \citep{goldowsky2023localizing,meng2022locating} to measure the Indirect Effect (IE) of each attention head $(l, h)$: the change in refusal score when the head's output from a harmless run is patched into a harmful run, computed separately for each task condition. This lets us compare each head's IE on task-wrapped inputs to its IE on direct harmful inputs, identifying universal heads (consistent IE across tasks) and task-specialized heads (high IE on direct harmful, low IE on task-wrapped).

\paragraph{Vocabulary projection.}
We apply the logit lens \citep{nostalgebraist2020logit} to track prediction entropy $H^{(\ell)}$ at each layer during steering interventions at varying magnitudes $\alpha$, and identify the commitment layer---the earliest layer at which refusal token probability exceeds 0.5 and remains stable through the final layer.

\paragraph{Experimental setup.}
We study LLaMA-3.1-8B-Instruct, LLaMA-2-7B-Chat, Gemma-7B-IT, and Qwen1.5-7B-Chat. For each model we construct per-task contrastive datasets across six instruction types (sentiment analysis, translation, summarization, QA, paraphrasing, and code generation), with 200 harmful and 200 harmless samples per task drawn from AdvBench \citep{zou2023advbench}, HarmBench \citep{mazeika2024harmbench}, JailbreakBench \citep{chao2024jailbreakbench}, and Alpaca \citep{taori2023alpaca} respectively. All circuit analysis uses TransformerLens \citep{nanda2022transformerlens}; refusal classification uses LLaMA-Guard-2 \citep{metallamaguard2024}.

%% -----------------------------------------------------------
\section{Results}
%% -----------------------------------------------------------
\label{sec:results}

\subsection{The Refusal Direction Is Task-Conditioned}

We first replicate the DIM finding: the universal $\hat{\mathbf{r}}$ achieves bypass rates between \ph{0.78} and \ph{0.91} across our four models, consistent with published results. Table~\ref{tab:tcrg} then shows \tcrg{} scores across models, all of which are non-trivial ($p < 0.001$, permutation test).

\begin{table}[h]
\centering
\small
\setlength{\tabcolsep}{6pt}
\begin{tabular}{lcc}
\toprule
Model & \tcrg{} & Max $\Delta_\text{cos}$ \\
\midrule
LLaMA-3.1-8B & \ph{0.38} & \ph{0.61} \\
LLaMA-2-7B   & \ph{0.31} & \ph{0.54} \\
Gemma-7B     & \ph{0.42} & \ph{0.67} \\
Qwen1.5-7B   & \ph{0.29} & \ph{0.49} \\
\bottomrule
\end{tabular}
\caption{\tcrg{} scores and maximum pairwise cosine distance between per-task refusal directions. All models show significant task-conditioned geometry ($p < 0.001$, permutation test). Largest divergences occur between translation and code generation.}
\label{tab:tcrg}
\end{table}

Figure~\ref{fig:bypass} shows the practical consequence: the universal direction achieves only \ph{0.41} bypass on sentiment and \ph{0.37} on translation, while task-specific directions achieve \ph{0.79} and \ph{0.81} on the same tasks. The gap is an order of magnitude larger than for direct harmful inputs (\ph{+0.04}).

\begin{figure}[t]
  \centering
  \fbox{
    \begin{minipage}{0.92\linewidth}
      \centering
      \vspace{0.5em}
      \textbf{Figure 2 --- Per-task bypass gap}

      \smallskip
      \emph{Suggested content:} Horizontal bar chart (one bar per task condition, ordered by bypass gap magnitude). Each bar is split into two segments: gray for universal $\hat{\mathbf{r}}$ bypass, colored extension for the additional gain from task-specific $\hat{\mathbf{r}}^{(k)}$. Sentiment and translation have the longest colored extensions (\ph{$+0.38$, $+0.44$}); harmful direct has almost none (\ph{$+0.04$}).
      \vspace{0.5em}
    \end{minipage}
  }
  \caption{Per-task bypass efficacy of the universal direction (gray) vs.\ task-specific direction (color). The bypass gap is largest for the most over-refused task types, directly reflecting the directional mismatch quantified by \tcrg{}.}
  \label{fig:bypass}
\end{figure}

\subsection{Two Classes of Refusal Heads and Neurons}

Activation patching reveals a clean dichotomy (Figure~\ref{fig:heads}). Universal heads (e.g., L14.H6, L18.H11; IE ratio $>$ 0.7) maintain strong causal influence across all task types and correspond well to the heads identified in prior circuit analyses \citep{arditi2024refusal}. Task-specialized heads (e.g., L16.H2, L12.H4; IE ratio $<$ 0.4) are strongly active for direct harmful inputs but weakly active for task-wrapped sentiment and translation.

At the neuron level, Table~\ref{tab:neurons} shows that sentiment and translation activate 43--51\% of their top refusal neurons from a task-unique population that does not overlap with the universal refusal neuron set. Ablating the universal DIM direction suppresses shared neurons but leaves task-unique neurons intact---a direct mechanistic account of why over-refusal persists after universal ablation.

\begin{figure}[t]
  \centering
  \fbox{
    \begin{minipage}{0.92\linewidth}
      \centering
      \vspace{0.5em}
      \textbf{Figure 3 --- Head IE profiles across tasks}

      \smallskip
      \emph{Suggested content:} Heatmap with attention heads on the y-axis (top-20 by mean IE, sorted by ratio) and six task conditions on the x-axis. Color encodes normalized IE. Top rows (universal heads) show uniformly warm color; bottom rows (task-specialized heads) show warm color only in the ``harmful direct'' column and cool elsewhere. A dashed horizontal rule separates the two groups.
      \vspace{0.5em}
    \end{minipage}
  }
  \caption{Indirect Effect (IE) of top-20 refusal-mediating attention heads across six task conditions. Universal heads maintain consistent causal influence; task-specialized heads are active only for direct harmful inputs.}
  \label{fig:heads}
\end{figure}

\begin{table}[h]
\centering
\small
\setlength{\tabcolsep}{5pt}
\begin{tabular}{lcc}
\toprule
Task & Universal overlap & Task-unique \\
\midrule
Harmful direct & 100\% & 0\% \\
Code           & 93\%  & 7\% \\
Summarization  & 83\%  & 17\% \\
QA             & 77\%  & 23\% \\
Sentiment      & 57\%  & 43\% \\
Translation    & 49\%  & 51\% \\
\bottomrule
\end{tabular}
\caption{Overlap between per-task top-50 refusal neurons and the universal refusal neuron set. Tasks with high task-unique fractions correspond exactly to the most over-refused conditions.}
\label{tab:neurons}
\end{table}

\subsection{Vocabulary Projection Reveals Pre-Failure Signals and Task-Dependent Commitment}

Figure~\ref{fig:entropy} shows prediction entropy at each layer during steering at varying magnitudes. Early-layer entropy ($\ell \leq 10$) rises sharply at $\alpha \geq 2.0$ for sentiment and translation inputs, predicting output incoherence before it is observable in the final generation. Crucially, this early-entropy signal appears at lower $\alpha$ values for task-wrapped inputs than for direct harmful inputs, reflecting the lower collapse thresholds of their higher-dimensional representations (Table~\ref{tab:thresholds}).

\begin{table}[h]
\centering
\small
\begin{tabular}{lcc}
\toprule
Task & Collapse $\alpha^*$ & Commit layer \\
\midrule
Code generation & \ph{3.1} & \ph{15.1} \\
Harmful direct  & \ph{2.8} & \ph{14.3} \\
Summarization   & \ph{2.2} & \ph{14.8} \\
Sentiment       & \ph{1.6} & \ph{11.4} \\
Translation     & \ph{1.4} & \ph{10.8} \\
\bottomrule
\end{tabular}
\caption{Per-task over-steering collapse threshold $\alpha^*$ and refusal commitment layer for LLaMA-3.1-8B. Both measures correlate with over-refusal severity: tasks that are most over-refused commit earliest and are most sensitive to high-magnitude interventions.}
\label{tab:thresholds}
\end{table}

The commitment layer analysis (Table~\ref{tab:thresholds}) shows that sentiment and translation commit to refusal near layer 11, roughly three layers before direct harmful inputs commit at layer 14. Steering applied at a fixed layer within this gap is post-commitment for task-wrapped inputs but pre-commitment for direct inputs---a mechanistic explanation for why the same intervention magnitude and layer produce very different outcomes across task types.

\begin{figure}[t]
  \centering
  \fbox{
    \begin{minipage}{0.92\linewidth}
      \centering
      \vspace{0.5em}
      \textbf{Figure 4 --- Vocabulary projection entropy and commitment curves}

      \smallskip
      \emph{Suggested content:} Two-panel figure.
      \textbf{Left:} Line plot with layer on the x-axis and prediction entropy (bits) on the y-axis. Separate lines for $\alpha \in \{0.5, 1.0, 2.0, 4.0, 8.0\}$ for sentiment inputs. Lines diverge sharply near $\ell = 8$ for $\alpha \geq 2.0$, marking the pre-failure region with a shaded band.
      \textbf{Right:} Commitment curves for five task conditions---probability of refusal token vs.\ layer. Sentiment and translation lines cross the 0.5 threshold near layer 11; harmful direct crosses near layer 14. The ``intervention window'' (pre-commitment region) is shaded per task.
      \vspace{0.5em}
    \end{minipage}
  }
  \caption{Left: Early-layer entropy rises predictably before output incoherence, providing a pre-failure diagnostic. Right: Commitment layer varies by task; over-refused task types commit 2--3 layers earlier, reducing the effective intervention window.}
  \label{fig:entropy}
\end{figure}

\subsection{Task-Conditioned Refusal Features Are Multi-Dimensional}

Following \citet{engels2025nonlinear}, we apply PCA to per-example refusal directions per task and measure intrinsic dimensionality (90\% variance threshold). Figure~\ref{fig:geometry} shows that direct harmful-input refusal is near one-dimensional (intrinsic dim $\approx$ \ph{2.1}), while sentiment and translation require \ph{4.8}--\ph{5.2} dimensions to capture 90\% of variance. Correspondingly, linear probes reach \ph{0.93} accuracy with 1D for harmful direct but only \ph{0.71} for sentiment; moving to 5D recovers \ph{$+0.18$} accuracy for task-conditioned conditions.

\begin{figure}[t]
  \centering
  \fbox{
    \begin{minipage}{0.92\linewidth}
      \centering
      \vspace{0.5em}
      \textbf{Figure 5 --- Feature geometry and probe accuracy}

      \smallskip
      \emph{Suggested content:} Two-panel figure.
      \textbf{Left:} PCA scree plot for per-task refusal directions (six tasks overlaid). Harmful direct (blue) drops sharply after component 1; sentiment (red) and translation (orange) remain elevated through component 5. The 90\% variance threshold is marked with a dashed line.
      \textbf{Right:} Grouped bar chart of probe accuracy for 1D, 2D, and 5D probes across the six task conditions. Direct harmful bars are nearly flat; sentiment and translation bars rise steeply, illustrating the multi-dimensional gain.
      \vspace{0.5em}
    \end{minipage}
  }
  \caption{Refusal feature geometry by task condition. Direct harmful-input refusal is near one-dimensional; task-conditioned over-refusal requires multi-dimensional representations, explaining why single-direction steering is insufficient.}
  \label{fig:geometry}
\end{figure}

This geometric structure provides a unified explanation for the lower collapse thresholds of task-wrapped inputs: large-magnitude additive interventions applied along a one-dimensional direction are orthogonal to many components of the higher-dimensional task-conditioned subspace, disrupting the representation as a whole rather than suppressing refusal specifically.

%% -----------------------------------------------------------
\section{Discussion}
%% -----------------------------------------------------------
\label{sec:discussion}

The four analyses above converge on a consistent mechanistic picture, illustrated in Figure~\ref{fig:overview}. In early layers, task identity is encoded as stable trajectory patterns in the residual stream. Mid-layer attention heads and MLP neurons read this task signal: shared universal heads activate across all conditions, while task-specialized heads activate specifically for task-wrapped harmful inputs. The resulting refusal representations are geometrically richer for task-conditioned conditions---occupying higher-dimensional subspaces---and commitment solidifies at task-dependent layers.

This picture connects directly to the empirical failure modes documented in prior work. The unreliability patterns identified by \citet{braun2025unreliability}---magnitude sensitivity, distributional distance, feature entanglement---are natural consequences of the circuit structure we document: task-unique neurons produce entanglement in mixed-task DIM extraction, geometric distance between universal and task-specific directions manifests as the distributional gap, and higher-dimensional task representations lower the collapse threshold. The abstention findings of \citet{hedstrom2025steer} correspond precisely to the post-commitment regime: their criterion for abstaining from steering is well-operationalized as detecting whether the intervention falls after the task-specific commitment layer. The per-category failure of universal truthfulness vectors \citep{wang2025adaptive} has an analogous explanation: different hallucination categories engage distinct circuit subpopulations with their own directional geometry, just as different task types do for refusal.

Our findings also clarify the scope of the DIM localization result \citep{arditi2024refusal}: the single-direction account is accurate for direct harmful inputs where the refusal feature is near one-dimensional and task-unique neuron populations are not engaged. It becomes an approximation when task context activates specialized sub-circuits. This is a refinement of scope rather than a contradiction.

Practically, the findings suggest three design principles for steering methods: extract directions separately per task type to minimize directional mismatch; use early-layer entropy as an adaptive signal to prevent over-steering collapse; and target interventions before the task-specific commitment layer rather than at a fixed layer.

%% -----------------------------------------------------------
\section{Conclusion}
%% -----------------------------------------------------------

We have provided a mechanistic account of why steering methods for refusal control succeed in some instruction contexts and fail in others. Through activation patching, vocabulary projection, and directional geometry analysis, we find that refusal is mediated by task-conditioned circuits that are geometrically heterogeneous: task-specialized attention heads and neurons engage selectively, refusal features span higher-dimensional subspaces for over-refused task types, and commitment to refusal solidifies at task-dependent layers. The \tcrg{} score quantifies this heterogeneity, and early-layer entropy provides a practical pre-failure diagnostic. We hope these findings contribute to the growing understanding of how safety behaviors are mechanistically structured, and suggest a productive direction for more reliable, context-aware steering methods.

%% -----------------------------------------------------------
\section*{Limitations}
%% -----------------------------------------------------------

Our analysis requires full activation access and is limited to open-source models. We study six task types; richer instruction categories (multi-turn, agentic, tool-use) remain to be examined. Circuit analysis operates within the residual stream linearity approximation. All numerical values are estimated placeholders pending full experimental execution.

%% -----------------------------------------------------------
\bibliography{custom}
%% -----------------------------------------------------------

\appendix

\section{Dataset Details}
\label{sec:appendix_data}

Task templates: (1) Sentiment---``Analyze the sentiment of the following: \{content\}''; (2) Translation---``Translate the following to French: \{content\}''; (3) Summarization---``Summarize the following: \{content\}''; (4) QA---``Answer questions about: \{content\}. Question: What is the main topic?''; (5) Code---``Write a Python function that: \{content\}''; (6) Paraphrasing---``Paraphrase the following: \{content\}''. Harmful texts: AdvBench \citep{zou2023advbench}, HarmBench \citep{mazeika2024harmbench}, JailbreakBench \citep{chao2024jailbreakbench}. Harmless texts: Alpaca \citep{taori2023alpaca}.

\section{TCRG Significance Testing}
\label{sec:appendix_stats}

TCRG significance is assessed via permutation test (10,000 permutations): task labels are randomly shuffled and TCRG recomputed each time to produce a null distribution. All $p$-values are two-tailed.

\section{Extended Circuit Results}
\label{sec:appendix_circuits}

\todo{Full head IE table for all models; per-layer neuron activation heatmaps showing task-specialized clustering in layers 12--18.}

\end{document}